Prior benchmarks of geospatial formats and engines record what has been measured, and on which datasets, across single-node spatial databases and Spark-based distributed engines \parencite{ray2011_jackpine_spatial_database_benchmark, pandey2018_how_good_modern_spatial_analytics}. A second line of work on benchmarking methodology for data systems prescribes how those measurements should be obtained and reported, covering controlled resource measurement and the statistical comparison of systems across repeated runs in variable cloud environments. Together, the two strands mark the gaps this thesis addresses. Spatial benchmarks rarely measure network transfer, and have not placed cloud-native and traditional formats and engines on one cloud workload with the statistical rigor the methodology literature recommends. These gaps motivate the research questions the thesis investigates, which set cloud-native formats and engines against traditional ones, and single-node against distributed execution, in performance and operational cost across shared spatial workloads.

\section{Geospatial format and engine benchmarks}
\label{related-work:geospatial-format-and-engine-benchmarks}

An early portable benchmark for spatial database engines is the Jackpine suite of \textcite{ray2011_jackpine_spatial_database_benchmark}, designed to fill the gap left by the absence of a TPC-Spatial standard. Jackpine combines micro queries based on the \acrshort{ogc} DE-9IM topological model with six macro workload scenarios, including map search and browsing, geocoding, flood-risk analysis, land information management, and toxic-spill analysis. The benchmark was applied to PostgreSQL, MySQL, and Informix on TIGER data for Texas, combined with City of Austin and Travis County GIS datasets, with the test harness running on the same machine as the database server. The benchmark targets single-node engines and predates the current cloud-native formats.

For distributed spatial systems, \textcite{pandey2018_how_good_modern_spatial_analytics} compared five Spark-based systems on Amazon \acrshort{emr} clusters of \num{1} to \num{16} nodes. The systems evaluated were SpatialSpark, GeoSpark, Magellan, Simba, and LocationSpark, run against \acrfull{osm} data comprising approximately \num{500} million records across point, line-string, rectangle, and polygon datasets. The study evaluated range, $k$-nearest-neighbor, distance-join, spatial-join, and $k$NN-join queries, reporting query throughput, total run-time, peak execution memory, and intra-cluster shuffle read and write volumes. GeoSpark, later renamed Apache Sedona, supported the broadest set of data types and queries and delivered the best performance in most cases, although it lacked $k$NN-join support. End-to-end network transfer between client and cluster was not measured.

Empirical comparisons of cloud-native vector formats against traditional formats appear mainly in community and vendor reports. \textcite{holmes2023_geoparquet_performance} compared GeoParquet, Shapefile, GeoPackage, and FlatGeobuf on the Google Open Buildings dataset, measuring write times and file sizes for inputs ranging from \num{498} megabytes to \num{21} gigabytes. GeoParquet produced the smallest files in every case, while Shapefile failed on the largest dataset. \textcite{flatgeobuf2024_flatgeobuf_benchmarks} reported relative read times for FlatGeobuf and Shapefile on Danish \acrshort{osm} roads (\num{906602} linestrings) and Danish cadastral polygons (\num{2511772} polygons), with FlatGeobuf reading the full datasets \num{2.2} and \num{8.3} times faster than Shapefile, respectively. Both studies use local disk storage and report only timing and file size; neither measures read or transfer over a network.

Two limitations recur across this literature. First, network transfer is rarely measured. \textcite{tang2020_locationspark_distributed_spatial_query} treats network communication cost as a system-design concern that motivates the spatial bitmap filter optimization rather than as a reported experimental metric, and \textcite{folkerts2012_benchmarking_in_the_cloud} argues that cloud benchmarks must report end-to-end performance jointly with pricing, of which data transfer is a major component. This thesis reports bytes received and bytes sent as primary metrics and derives the corresponding cost, addressing both axes directly. Second, single-node and distributed engines are rarely compared in the same study; \textcite{pandey2018_how_good_modern_spatial_analytics} compares only Spark-based systems, while Jackpine and similar single-node benchmarks predate the Spark-based spatial systems that emerged in the mid-2010s.
\section{Benchmarking methodology for data systems}
\label{related-work:benchmarking-methodology-for-data-systems}

\textcite{beyer_2015_benchmarking_and_resource_measurement} identifies a set of requirements for reliable benchmarking, including accurate measurement of CPU time and memory, reliable termination of processes, and a controlled execution environment. They show that classical Linux mechanisms such as \texttt{time} and \texttt{ulimit} fail to meet these requirements, primarily because they do not account for child processes correctly, and propose a framework based on Linux control groups (cgroups), a kernel feature that tracks and limits resource usage for groups of processes as a unit. The \acrfull{jcgm} supplies the underlying terminology used in this thesis to describe measurement quality, including the distinction between accuracy, precision, repeatability, and reproducibility \parencite{jcgm2012_international_vocabulary_of_metrology}.

Timing-measurement pitfalls have received particular attention. \textcite{chen2016_robust_benchmarking_noisy_environments} showed that consecutive run-time measurements are not necessarily independent and identically distributed, and recommended the minimum estimator over repeated short runs as a robust summary. \textcite{hoefler2015_scientific_benchmarking_parallel} reviewed 120 papers from three top conferences in high-performance computing and codified twelve rules for reporting parallel-system performance, emphasizing that distributions, not point averages, should be reported. \textcite{raasveldt2018_fair_benchmarking_database_systems} cataloged pitfalls specific to database benchmarking, including cold-versus-hot caches, untuned baselines, comparison of systems with different functionality, and ignored preprocessing and index-creation time, each of which can flip a published comparison.

Cloud environments add a further source of variation. \textcite{iosup2011_performance_variability_cloud_services} analyzed ten production services on \acrlong{aws} and Google App Engine over a year and found yearly, monthly, weekly, and daily patterns in service performance, with coefficients of variation above $1.10$ indicating high variability for several services. \textcite{schad2010_runtime_measurements_cloud} demonstrated that nominally identical \acrfull{ec2} instances differ measurably in \acrshort{cpu}, \acrshort{io}, and network throughput, with performance falling into two distinct bands tied to the underlying virtual hardware. \textcite{leitner2016_patterns_chaos_iaas} extended these results across four \acrshort{iaas} providers, including \acrshort{aws} \acrshort{ec2}, Google Compute Engine, Microsoft Azure, and IBM SoftLayer, codifying fifteen hypotheses about cloud performance variability and reporting that \acrshort{ec2} and Azure were substantially less predictable than the other two providers at the time of their study. \textcite{laaber2019_software_microbenchmarking_cloud} measured $4.5$ million microbenchmark runs across three public-cloud providers and found that running test and control experiments on the same instance in randomized order, combined with the Wilcoxon rank-sum test and bootstrapped confidence intervals, reliably detects slowdowns of $10\%$ or less.

Statistical comparison of systems is rarely reported with effect sizes. \textcite{arcuri2014_hitchhiker_statistical_tests} recommend reporting the Vargha--Delaney $\hat{A}_{12}$ effect size alongside non-parametric significance tests, with values near $0.5$ interpreted as no effect, $\geq 0.56$ as a small effect, $\geq 0.64$ as a medium effect, and $\geq 0.71$ as a large effect. The methodology chapter follows this approach.


\section{Research gaps}
\label{related-work:research-gaps}
Three gaps in the literature surveyed above motivate the research questions of this thesis.

First, network transfer accounting is absent from almost all spatial benchmarks, despite being a primary driver of cost and latency in cloud delivery. \textcite{tang2020_locationspark_distributed_spatial_query} treats network communication as a system-design concern rather than a reported experimental metric, and \textcite{folkerts2012_benchmarking_in_the_cloud} argues that cloud benchmarks must report end-to-end performance jointly with pricing, of which data transfer is a major component. The link from format internals to client-observed cost has therefore not been measured directly in the spatial benchmarking literature.

Second, no peer-reviewed benchmark compares vector cloud-native formats (GeoParquet, PMTiles) against traditional formats and engines (Shapefile, PostGIS, WMS-based vector tiles) on the same workload. The empirical comparisons that exist appear in community and vendor reports \parencite{holmes2023_geoparquet_performance, flatgeobuf2024_flatgeobuf_benchmarks}, are restricted to local-disk timing and file size, and do not place cloud-native and traditional formats side by side on cloud storage.

Third, no published benchmark compares single-node engines (PostGIS, DuckDB) against distributed engines (Apache Sedona on Databricks) on the same vector workload. \textcite{pandey2018_how_good_modern_spatial_analytics} compares only Spark-based systems, while Jackpine \parencite{ray2011_jackpine_spatial_database_benchmark} and similar single-node benchmarks predate the Spark-based spatial systems that emerged in the mid-2010s.

The first two gaps motivate \acrshort{rq}1, which compares cloud-native formats and engines against traditional ones in performance and operational cost across analytical and visualization spatial workloads. The third gap motivates \acrshort{rq}2, which extends the comparison to horizontal scaling of distributed cloud-native engines against single-node engines on the same workloads. The methodological gap discussed in Section~\ref{ch:methodology}, the absence of effect-size reporting and uncertainty quantification in geospatial benchmarks, is addressed as a methodological contribution rather than a separate research question.
\section{Research questions}
\label{related-work:research-questions}
This thesis addresses three research questions, each following from the gaps identified above.

The first concerns single-machine query performance. \textbf{RQ1:} How does the cloud-native combination of GeoParquet and DuckDB compare against the traditional alternatives PostGIS and Shapefile in wall-clock query time, network transfer, and operational cost across spatial intersection predicates, distance-based nearest-neighbor search, and bounding-box range filtering on Norwegian building data at small and large dataset tiers?

The second turns to distributed execution and scaling. \textbf{RQ2:} For a national-scale spatial join on Norwegian building data, how do broadcast-join and partitioned-join execution strategies on Apache Sedona compare to single-node PostGIS and DuckDB in wall-clock time, scaling behavior, network transfer, and operational cost across small, medium, and large dataset tiers?

The third examines whether a single winner holds across the comparison. \textbf{RQ3:} Does the best-performing format-and-engine combination in RQ1 remain consistent, or do winners differ across query patterns and outcome dimensions?
